\section{Resumen ejecutivo}

PayCore es un motor transaccional para una billetera digital con saldo
simulado. Su objetivo no es parecerse a una aplicación comercial, sino
demostrar que las garantías que un sistema de pagos debe ofrecer se pueden
construir, probar y defender con argumentos verificables.

El sistema implementa un libro mayor de doble entrada inmutable como fuente de
verdad del saldo, control de concurrencia respaldado por la base de datos,
idempotencia persistente confirmada en la misma transacción que el movimiento
de dinero, autenticación con Argon2id y JWT, y una interfaz web que deja claro
en todo momento que los fondos no son reales.

\subsection*{Qué se construyó}

Un monolito modular en Python con FastAPI y SQLAlchemy sobre PostgreSQL~16,
una interfaz en React con TypeScript, 96 pruebas automatizadas que se ejecutan
contra un PostgreSQL real, y un conjunto de mediciones de carga con Locust
cuyos datos crudos están versionados en el repositorio.

\subsection*{Resultados verificados}

\begin{itemize}[leftmargin=1.2em]
  \item Las seis pruebas funcionales exigidas (transferencia correcta, fondos
    insuficientes, solicitud duplicada, cien solicitudes concurrentes, fallos
    antes y después del \codigo{COMMIT}, y seguridad) pasan y se comprueban
    \textbf{consultando la base de datos}, no las respuestas HTTP.
  \item Bajo contención máxima, con cien usuarios durante sesenta segundos y
    cuatro procesos de API compitiendo, el libro mayor sumó exactamente cero,
    ninguna cuenta quedó descuadrada y ningún saldo resultó negativo, con las
    dos estrategias de concurrencia implementadas.
  \item La estrategia pesimista confirmó 148,3 transferencias por segundo
    frente a 10,4 de la estrategia \codigo{SERIALIZABLE}, cuya tasa de fallo
    alcanzó el 83,72\,\% por agotamiento de reintentos. La diferencia solo es
    visible contando filas en PostgreSQL: en peticiones por segundo la brecha
    aparentaba ser de apenas un 57\,\%.
\end{itemize}

\subsection*{Dos defectos reales encontrados por las pruebas}

El valor de este trabajo no está solo en lo que funciona, sino en lo que falló
y fue detectado:

\begin{enumerate}[leftmargin=1.4em]
  \item Una \textbf{actualización perdida} que ninguna prueba de un solo hilo
    podía ver. El bloqueo \codigo{SELECT FOR UPDATE} se adquiría
    correctamente, pero el mapa de identidad de SQLAlchemy devolvía el objeto
    ya cargado sin refrescar sus atributos, de modo que la comprobación de
    fondos evaluaba un saldo obsoleto. Con el defecto presente, la prueba de
    concurrencia aprobaba 88 transferencias de 100 cuando el máximo posible
    eran 3, y las 76 pruebas existentes pasaban todas.
  \item Un \textbf{agotamiento del pool de conexiones} que la carga destapó.
    La sesión de base de datos se obtenía como dependencia del framework, lo
    que retenía la conexión fuera del hilo que la usaba. Corregirlo multiplicó
    por 7,4 el rendimiento de lectura y redujo el peor tiempo de respuesta de
    30\,664\,ms a 1\,100\,ms.
\end{enumerate}

\subsection*{Lo que PayCore no es}

No es una plataforma bancaria de producción. Funciona sobre un solo nodo de
base de datos sin alta disponibilidad, su limitación de peticiones cuenta en
memoria de cada proceso, sus tokens no admiten revocación inmediata y su
despliegue con Docker Compose no es un despliegue real. Estas limitaciones
están enumeradas en la sección~\ref{sec:limitaciones} y declaradas también en
el repositorio, porque conocerlas es parte del trabajo.
